\documentclass[11pt]{article}

% Change "review" to "final" to generate the final (sometimes called camera-ready) version.
% Change to "preprint" to generate a non-anonymous version with page numbers.
\usepackage[]{acl}
\usepackage{times}
\usepackage{latexsym}

% For proper rendering and hyphenation of words containing Latin characters (including in bib files)
\usepackage[T1]{fontenc}
% For Vietnamese characters
% \usepackage[T5]{fontenc}
% See https://www.latex-project.org/help/documentation/encguide.pdf for other character sets

% This assumes your files are encoded as UTF8
\usepackage[utf8]{inputenc}
\usepackage{microtype}
\usepackage{inconsolata}
\usepackage{graphicx}

% If the title and author information does not fit in the area allocated, uncomment the following
%
%\setlength\titlebox{<dim>}
%
% and set <dim> to something 5cm or larger.

\title{Group X Progress Report:\\My Group's Project Name}


\author{First Author, Second Author, Third Author \\
  \texttt{\{macid1,macid2,macid3\}@mcmaster.ca} }

%\author{
%  \textbf{First Author\textsuperscript{1}},
%  \textbf{Second Author\textsuperscript{1,2}},
%  \textbf{Third T. Author\textsuperscript{1}},
%  \textbf{Fourth Author\textsuperscript{1}},
%\\
%  \textbf{Fifth Author\textsuperscript{1,2}},
%  \textbf{Sixth Author\textsuperscript{1}},
%  \textbf{Seventh Author\textsuperscript{1}},
%  \textbf{Eighth Author \textsuperscript{1,2,3,4}},
%\\
%  \textbf{Ninth Author\textsuperscript{1}},
%  \textbf{Tenth Author\textsuperscript{1}},
%  \textbf{Eleventh E. Author\textsuperscript{1,2,3,4,5}},
%  \textbf{Twelfth Author\textsuperscript{1}},
%\\
%  \textbf{Thirteenth Author\textsuperscript{3}},
%  \textbf{Fourteenth F. Author\textsuperscript{2,4}},
%  \textbf{Fifteenth Author\textsuperscript{1}},
%  \textbf{Sixteenth Author\textsuperscript{1}},
%\\
%  \textbf{Seventeenth S. Author\textsuperscript{4,5}},
%  \textbf{Eighteenth Author\textsuperscript{3,4}},
%  \textbf{Nineteenth N. Author\textsuperscript{2,5}},
%  \textbf{Twentieth Author\textsuperscript{1}}
%\\
%\\
%  \textsuperscript{1}Affiliation 1,
%  \textsuperscript{2}Affiliation 2,
%  \textsuperscript{3}Affiliation 3,
%  \textsuperscript{4}Affiliation 4,
%  \textsuperscript{5}Affiliation 5
%\\
%  \small{
%    \textbf{Correspondence:} \href{mailto:email@domain}{email@domain}
%  }
%}

\begin{document}
\maketitle
% \begin{abstract}
% \end{abstract}

\section{Introduction (DONE)}

%Here, write a brief introduction to the problem you are solving. This can be adapted from your problem description and motivation from the original proposal. This should be around 0.25-0.5 pages.

We used an existing algorithm to play 2048 and collected the board state before each recorded move and the move it chose as training data. We trained the model on these data without explicitly giving it the algorithm's rules. We aim to find out whether the model can learn strategies and imitate the algorithm using only examples of board states and chosen moves.

\section{Related Work}

Here, talk about the related work you encountered for your approach. Cite at least 5 references. Refer to item 2. No one has done exactly your task? Write about the most similar thing you can find. This should be around 0.25-0.5 pages.

\section{Dataset (DONE)}

%You should write about your dataset here, following the guidelines regarding item 1. This section may be 0.5-1 pages. Depending on your specific dataset, you may want to include subsections for the preprocessing, annotation, etc.

Our dataset was collected using a scraping technique to generate games from the front-end implementation of the existing 2048 ML model our program utilizes: https://game-difficulty.github.io/2048EndgameTablebase/. Using Selenium and Beautifulsoup libraries, our team created a program that would open the 2048 simulator in a browser window and increment through games one turn at a time, saving the board state and current score on each move. The boardstate was collected from the corresponding hexadecimal code given at the top of the page. This code was split into individual board tiles (1 digit for each tile), and each tile was converted into its corresponding decimal value. Each board tile was stored in the logarithmic form (e.g. 4 would be stored as 2, since 2^2 = 4). Each board tiles name was adherent to the A-D:1-4 formatting, with the letter representing the row, and the number representing the column. At the end of a game, the board would be cleared and reset to a valid starting state that was randomly generated by the program. The dataset also includes an estimate for what move was made, based off of the board state before and after each move was made. When a valid estimate for the previous move could not be made, the program defaulted to a 0 value for the move value. These zeroes were filtered out in the preprocessing stage. The program exports collected data to an excel sheet, and multiple datasheets were created with varying amounts of games and turns. 

\section{Features (DONE)}

%Describe any features you used for your model, or how your data was input to your model. Are you doing feature engineering or feature selection? Are you learning embeddings? Is it all part of one neural network? Refer to item 2. This may range from 0.25 pages to 0.5 pages.

We represented each $4 \times 4$ board as 16 numbers, one for each cell, in row order. An empty cell was encoded as 0, and a tile with value $2^n$ was encoded as $n$. We kept all 16 features and did not use data augmentation. We fitted \texttt{StandardScaler} on the training data and used it to standardize both the training and test data before running k-NN.

\section{Implementation (DONE)}

%Describe your model and implementation here. Refer to item 4. This may take around a page.

We used scikit-learn to implement a k-nearest neighbors (k-NN)classifier with $k=5$. The model takes the 16 encoded tile values from a board and predicts one of four moves: up, down, left, or right. We calculated the scaling statistics from the training boards and used them to standardize both the training and test boards.

For each test board, the classifier finds the five closest training boards using Euclidean distance. Each neighbor gets one vote, and the move with the most votes is the prediction. Our k-NN model stores the training boards and predicts by voting among the five closest ones, so it does not need a loss function during training.

For the baseline, we ranked the four moves by how often they appeared in the training data. For each test board, the baseline checks the moves in that order and picks the first one that slides or merges a tile.

\section{Results and Evaluation (DONE)}

%How are you evaluating your model? What results do you have so far? What are your baselines? Refer to item 5. This may take around 0.5 pages.

We split the data by game, using the first 80\% of games for training and the remaining 20\% for testing. We did not use a separate validation set or cross-validation. Each game contributed 50 recorded moves. We measured test accuracy as the percentage of predictions that matched the original AI's moves. For each test board, the baseline chose the legal move that appeared most often in the training data.

As shown in Table~\ref{tab:move-accuracy}, k-NN beat the baseline in four of the five experiments. In the 200-game experiment, k-NN achieved 52.0\% accuracy, compared with 47.7\% for the baseline.

\begin{table}[htbp]
\centering
\begin{tabular}{rrr}
\hline
Games & k-NN($k=5$) & Baseline \\
\hline
5   & 48.0\% & 66.0\% \\
10  & 49.0\% & 41.0\% \\
50  & 54.8\% & 45.4\% \\
135 & 52.4\% & 51.0\% \\
200 & 52.0\% & 47.7\% \\
\hline
\end{tabular}
\caption{Move prediction accuracy by dataset size}
\label{tab:move-accuracy}
\end{table}

\section{Feedback and Plans}

Write about your plans for the remainder of the project. This should include a discussion of the feedback you received from your TA, and how you plan to improve your approach. Reflect on your implementation and areas for improvement. Refer to item 6. This may be around 0.5 pages.

\section{Template Notes}

You can remove this section or comment it out, as it only contains instructions for how to use this template. You may use subsections in your document as you find appropriate.

\subsection{Tables and figures}

See Table~\ref{citation-guide} for an example of a table and its caption.
See Figure~\ref{fig:experiments} for an example of a figure and its caption.


\begin{figure}[t]
  \includegraphics[width=\columnwidth]{example-image-golden}
  \caption{A figure with a caption that runs for more than one line.
    Example image is usually available through the \texttt{mwe} package
    without even mentioning it in the preamble.}
  \label{fig:experiments}
\end{figure}

\begin{figure*}[t]
  \includegraphics[width=0.48\linewidth]{example-image-a} \hfill
  \includegraphics[width=0.48\linewidth]{example-image-b}
  \caption {A minimal working example to demonstrate how to place
    two images side-by-side.}
\end{figure*}


\subsection{Citations}

\begin{table*}
  \centering
  \begin{tabular}{lll}
    \hline
    \textbf{Output}           & \textbf{natbib command} & \textbf{ACL only command} \\
    \hline
    \citep{Gusfield:97}       & \verb|\citep|           &                           \\
    \citealp{Gusfield:97}     & \verb|\citealp|         &                           \\
    \citet{Gusfield:97}       & \verb|\citet|           &                           \\
    \citeyearpar{Gusfield:97} & \verb|\citeyearpar|     &                           \\
    \citeposs{Gusfield:97}    &                         & \verb|\citeposs|          \\
    \hline
  \end{tabular}
  \caption{\label{citation-guide}
    Citation commands supported by the style file.
  }
\end{table*}

Table~\ref{citation-guide} shows the syntax supported by the style files.
We encourage you to use the natbib styles.
You can use the command \verb|\citet| (cite in text) to get ``author (year)'' citations, like this citation to a paper by \citet{Gusfield:97}.
You can use the command \verb|\citep| (cite in parentheses) to get ``(author, year)'' citations \citep{Gusfield:97}.
You can use the command \verb|\citealp| (alternative cite without parentheses) to get ``author, year'' citations, which is useful for using citations within parentheses (e.g. \citealp{Gusfield:97}).

\subsection{References}

\nocite{Ando2005,andrew2007scalable,rasooli-tetrault-2015}

Many websites where you can find academic papers also allow you to export a bib file for citation or bib formatted entry. Copy this into the \texttt{custom.bib} and you will be able to cite the paper in the \LaTeX{}. You can remove the example entries.

\subsection{Equations}

An example equation is shown below:
\begin{equation}
  \label{eq:example}
  A = \pi r^2
\end{equation}

Labels for equation numbers, sections, subsections, figures and tables
are all defined with the \verb|\label{label}| command and cross references
to them are made with the \verb|\ref{label}| command.
This an example cross-reference to Equation~\ref{eq:example}. You can also write equations inline, like this: $A=\pi r^2$.


% \section*{Limitations}

\section*{Team Contributions}

Write in this section a few sentences describing the contributions of each team member. What did each member work on? Refer to item 7.

% Bibliography entries for the entire Anthology, followed by custom entries
%\bibliography{custom,anthology-overleaf-1,anthology-overleaf-2}

% Custom bibliography entries only
\bibliography{custom}

% \appendix

% \section{Example Appendix}
% \label{sec:appendix}

% This is an appendix.

\end{document}
